% Source: physical PDF pp. 275--276 (printed pp. 252--253). Coverage: Chapter 20 opening, from the chapter heading through immediately before Section 20.1; one unnumbered display, citations [1]--[2], no figures or footnotes. Uncertainty: none.

We often find ourselves needing to know how an amplitude behaves when the
four-momentum brought in by one operator and out by another goes to infinity,
with all other external lines held at fixed four-momenta. For instance, we
will see in Section 20.6 that the total cross section for the scattering of an
electron by an initial hadron $H$, with arbitrary hadrons in the final state,
is given by unitarity as a linear combination (with known coefficients) of
the components of the amplitude
\[
 \int d^4x\,e^{-ik\cdot x}
 \bra{H}J^\mu(x)J^\nu(0)\ket{H},
\]
where $k$ is the four-momentum transferred from the electron to the hadrons,
and $J^\mu(x)$ is the electromagnetic current. In the case of deep inelastic
electric scattering the momentum $k$ that is carried in by one current
operator and out by the other is allowed to go to infinity. Similarly, in
studying the high momentum limit of various propagators and deriving
corresponding spectral function sum rules in Section 20.5, we shall encounter
the high momentum limit of similar Fourier transforms, but with the
one-hadron state $\ket{H}$ replaced with the vacuum.

If an operator product such as $J^\mu(x)J^\nu(0)$ were analytic in $x^\mu$,
then its Fourier transform would decrease exponentially as the Fourier
variable $k$ goes to infinity. The leading terms in the high momentum limit
of the Fourier transform arise from the singularities of the operator product
as the spacetime arguments approach one another. The study of such operator
products was initiated in 1969 by Wilson~\hyperref[ch20-ref-1]{[1]},
originally as an attempt to formulate a substitute for conventional quantum
field theory. As has happened earlier (for instance with dispersion
relations and Feynman's diagrammatic rules) the effort to bypass quantum
field theory led to valuable general results, but results that can best be
understood as general properties of quantum field theory.

The operator product expansion will be stated in Section 20.1. The standard
proof of this expansion was given in perturbation theory in 1970 by
Zimmerman~\hyperref[ch20-ref-2]{[2]}. In Section 20.1 we shall offer a
non-perturbative and simpler though less rigorous derivation based on the
path-integral formulation of field theory. Section 20.2 will present a
different perspective on the operator product expansion, in terms of the flow
of large momenta through Feynman graphs, which will lead us to a perturbative
proof.

There are several aspects of the operator product expansion that make it
particularly useful for drawing consequences from theories like quantum
chromodynamics. One such property, discussed in Section 20.3, is that the
functions describing the singularities in this expansion have a momentum
dependence governed by renormalization group expansions, so that in
asymptotically free theories they can be calculated at large momenta using
perturbation theory. Another aspect, shown in Section 20.4, is that these
functions exhibit the full symmetry of the underlying theory, unaffected by
possible spontaneous symmetry breaking. Applications will be considered in
Sections 20.5 and 20.6.
